% GENERATED FILE. Edit canonical lessons, then run npm run generate:books.
% canonical-source-hash: fnv1a64:c1850e124a94053d
% canonical-lessons: JA-C499-yuuenchi, JA-C499-doubutsuen, JA-C499-suizokukan, JA-C499-ensoku, JA-C499-kouraku

\chapter{Days Out}
\label{ch:ja-days-out}
\hlchaptermodality{499}

\begin{chapteropening}
\textbf{By the end of this chapter:} \emph{I can say where we go on a day out.}

Complete the last lesson of chapter 499: I can say where we go on a day out.
\end{chapteropening}

\section[\texorpdfstring{\ja{ゆうえんち} (yūenchi)}{ゆうえんち (yūenchi)}]{\ja{ゆうえんち} (yūenchi) — an amusement park}

\label{lesson:JA-C499-yuuenchi}



\begin{quote}
Before the new one: say the Japanese for nodding off, then the Japanese for paid leave.
\end{quote}

\subsection*{You'll want to know: \ja{ゆうえんち}}
\textbf{\ja{ゆうえんち}} — \emph{yūenchi} — \textquotedblleft{}an amusement park\textquotedblright{}.

\begin{quote}
\textbf{\ja{にちようびに} \ja{ゆうえんちへ} \ja{いきました。}} — \emph{nichiyōbi ni yūenchi e ikimashita} — \textquotedblleft{}On Sunday we went to the amusement park.\textquotedblright{}
\end{quote}

\begin{cousinweb}
\textbf{\ja{ゆうえんち}} (\emph{yūenchi}) is Sino-Japanese (kango): three Chinese characters read with their on'yomi, readings borrowed from Chinese. They are \emph{yū}, \textquotedblleft{}play\textquotedblright{}, \emph{en}, \textquotedblleft{}park\textquotedblright{} (Middle Chinese \emph{hjwon}) and \emph{chi}, \textquotedblleft{}ground\textquotedblright{} (\emph{dijH}) (KANJIDIC2; Middle Chinese in Baxter's transcription).
\end{cousinweb}

\begin{grammarlens}[title={What you've built}]
One more word for describing what you lived and felt.
\end{grammarlens}

\subsection*{Your turn}
\begin{itemize}
\raggedright
  \item \emph{Say it:} \emph{yūenchi}
  \item \emph{Say it:} \emph{yūenchi}, once more
  \item \emph{Say it:} \emph{nichiyōbi ni yūenchi e ikimashita} — On Sunday we went to the amusement park.
  \item \emph{Recall:} say the Japanese for nodding off, then the Japanese for paid leave, then say \emph{yūenchi} again
\end{itemize}

\begin{tcolorbox}[breakable,colback=teal!4,colframe=teal!35!black,title={Before you move on}]
What does \ja{ゆうえんち} mean? (\textquotedblleft{}An amusement park\textquotedblright{}.) Say it once more.
\end{tcolorbox}
\section[\texorpdfstring{\ja{どうぶつえん} (dōbutsuen)}{どうぶつえん (dōbutsuen)}]{\ja{どうぶつえん} (dōbutsuen) — a zoo}

\label{lesson:JA-C499-doubutsuen}



\begin{quote}
Before the new one: say the Japanese for paid leave, then the Japanese for an amusement park.
\end{quote}

\subsection*{You'll want to know: \ja{どうぶつえん}}
\textbf{\ja{どうぶつえん}} — \emph{dōbutsuen} — \textquotedblleft{}a zoo\textquotedblright{}.

\begin{quote}
\textbf{\ja{どうぶつえんで} \ja{さるを} \ja{みました。}} — \emph{dōbutsuen de saru o mimashita} — \textquotedblleft{}We saw monkeys at the zoo.\textquotedblright{}
\end{quote}

\begin{cousinweb}
\textbf{\ja{どうぶつえん}} (\emph{dōbutsuen}) is Sino-Japanese (kango): three Chinese characters read with their on'yomi, readings borrowed from Chinese. They are \emph{dō}, \textquotedblleft{}move\textquotedblright{} (Middle Chinese \emph{duwngX}), \emph{butsu}, \textquotedblleft{}thing\textquotedblright{} (\emph{mjut}) and \emph{en}, \textquotedblleft{}garden\textquotedblright{} (\emph{hjwon}) (KANJIDIC2; Middle Chinese in Baxter's transcription).
\end{cousinweb}

\begin{grammarlens}[title={What you've built}]
One more word for describing what you lived and felt.
\end{grammarlens}

\subsection*{Your turn}
\begin{itemize}
\raggedright
  \item \emph{Say it:} \emph{dōbutsuen}
  \item \emph{Say it:} \emph{dōbutsuen}, once more
  \item \emph{Say it:} \emph{dōbutsuen de saru o mimashita} — We saw monkeys at the zoo.
  \item \emph{Recall:} say the Japanese for paid leave, then the Japanese for an amusement park, then say \emph{dōbutsuen} again
\end{itemize}

\begin{tcolorbox}[breakable,colback=teal!4,colframe=teal!35!black,title={Before you move on}]
What does \ja{どうぶつえん} mean? (\textquotedblleft{}A zoo\textquotedblright{}.) Say it once more.
\end{tcolorbox}
\section[\texorpdfstring{\ja{すいぞくかん} (suizokukan)}{すいぞくかん (suizokukan)}]{\ja{すいぞくかん} (suizokukan) — an aquarium}

\label{lesson:JA-C499-suizokukan}



\begin{quote}
Before the new one: say the Japanese for an amusement park, then the Japanese for a zoo.
\end{quote}

\subsection*{You'll want to know: \ja{すいぞくかん}}
\textbf{\ja{すいぞくかん}} — \emph{suizokukan} — \textquotedblleft{}an aquarium\textquotedblright{}.

\begin{quote}
\textbf{\ja{すいぞくかんには} \ja{さめが} \ja{います。}} — \emph{suizokukan ni wa same ga imasu} — \textquotedblleft{}There are sharks at the aquarium.\textquotedblright{}
\end{quote}

\begin{cousinweb}
\textbf{\ja{すいぞくかん}} (\emph{suizokukan}) is Sino-Japanese (kango): three Chinese characters read with their on'yomi, readings borrowed from Chinese. They are \emph{sui}, \textquotedblleft{}water\textquotedblright{} (Middle Chinese \emph{sywijX}), \emph{zoku}, \textquotedblleft{}family\textquotedblright{} (\emph{dzuwk}) and \emph{kan}, \textquotedblleft{}building\textquotedblright{} (\emph{kwanH}) (KANJIDIC2; Middle Chinese in Baxter's transcription).
\end{cousinweb}

\begin{grammarlens}[title={What you've built}]
One more word for describing what you lived and felt.
\end{grammarlens}

\subsection*{Your turn}
\begin{itemize}
\raggedright
  \item \emph{Say it:} \emph{suizokukan}
  \item \emph{Say it:} \emph{suizokukan}, once more
  \item \emph{Say it:} \emph{suizokukan ni wa same ga imasu} — There are sharks at the aquarium.
  \item \emph{Recall:} say the Japanese for an amusement park, then the Japanese for a zoo, then say \emph{suizokukan} again
\end{itemize}

\begin{tcolorbox}[breakable,colback=teal!4,colframe=teal!35!black,title={Before you move on}]
What does \ja{すいぞくかん} mean? (\textquotedblleft{}An aquarium\textquotedblright{}.) Say it once more.
\end{tcolorbox}
\section[\texorpdfstring{\ja{えんそく} (ensoku)}{えんそく (ensoku)}]{\ja{えんそく} (ensoku) — a school outing; an excursion}

\label{lesson:JA-C499-ensoku}



\begin{quote}
Before the new one: say the Japanese for a zoo, then the Japanese for an aquarium.
\end{quote}

\subsection*{You'll want to know: \ja{えんそく}}
\textbf{\ja{えんそく}} — \emph{ensoku} — \textquotedblleft{}a school outing; an excursion\textquotedblright{}.

\begin{quote}
\textbf{\ja{らいしゅう} \ja{がっこうの} \ja{えんそくが} \ja{あります。}} — \emph{raishū gakkō no ensoku ga arimasu} — \textquotedblleft{}Next week there is a school outing.\textquotedblright{}
\end{quote}

\begin{cousinweb}
\textbf{\ja{えんそく}} (\emph{ensoku}) is Sino-Japanese (kango): two Chinese characters read with their on'yomi, readings borrowed from Chinese. They are \emph{en}, \textquotedblleft{}distant\textquotedblright{} (Middle Chinese \emph{hjwonX}) and \emph{soku}, \textquotedblleft{}leg\textquotedblright{} (\emph{tsjowk}) (KANJIDIC2; Middle Chinese in Baxter's transcription).
\end{cousinweb}

\begin{grammarlens}[title={What you've built}]
One more word for describing what you lived and felt.
\end{grammarlens}

\subsection*{Your turn}
\begin{itemize}
\raggedright
  \item \emph{Say it:} \emph{ensoku}
  \item \emph{Say it:} \emph{ensoku}, once more
  \item \emph{Say it:} \emph{raishū gakkō no ensoku ga arimasu} — Next week there is a school outing.
  \item \emph{Recall:} say the Japanese for a zoo, then the Japanese for an aquarium, then say \emph{ensoku} again
\end{itemize}

\begin{tcolorbox}[breakable,colback=teal!4,colframe=teal!35!black,title={Before you move on}]
What does \ja{えんそく} mean? (\textquotedblleft{}A school outing; an excursion\textquotedblright{}.) Say it once more.
\end{tcolorbox}
\section[\texorpdfstring{\ja{こうらく} (kōraku)}{こうらく (kōraku)}]{\ja{こうらく} (kōraku) — a pleasure trip; an outing}

\label{lesson:JA-C499-kouraku}



\begin{quote}
Before the new one: say the Japanese for an aquarium, then the Japanese for a school outing.
\end{quote}

\subsection*{You'll want to know: \ja{こうらく}}
\textbf{\ja{こうらく}} — \emph{kōraku} — \textquotedblleft{}a pleasure trip; an outing\textquotedblright{}.

\begin{quote}
\textbf{\ja{あきは} \ja{こうらくの} \ja{きせつです。}} — \emph{aki wa kōraku no kisetsu desu} — \textquotedblleft{}Autumn is the season for outings.\textquotedblright{}
\end{quote}

\begin{cousinweb}
\textbf{\ja{こうらく}} (\emph{kōraku}) is Sino-Japanese (kango): two Chinese characters read with their on'yomi, readings borrowed from Chinese. They are \emph{kō}, \textquotedblleft{}going\textquotedblright{} (Middle Chinese \emph{haeng}) and \emph{raku}, \textquotedblleft{}comfort\textquotedblright{} (KANJIDIC2; Middle Chinese in Baxter's transcription).
\end{cousinweb}

\begin{grammarlens}[title={What you've built}]
One more word for describing what you lived and felt.
\end{grammarlens}

\subsection*{Your turn}
\begin{itemize}
\raggedright
  \item \emph{Say it:} \emph{kōraku}
  \item \emph{Say it:} \emph{kōraku}, once more
  \item \emph{Say it:} \emph{aki wa kōraku no kisetsu desu} — Autumn is the season for outings.
  \item \emph{Recall:} say the Japanese for an aquarium, then the Japanese for a school outing, then say \emph{kōraku} again
\end{itemize}

\begin{tcolorbox}[breakable,colback=teal!4,colframe=teal!35!black,title={Before you move on}]
What does \ja{こうらく} mean? (\textquotedblleft{}A pleasure trip; an outing\textquotedblright{}.) Say it once more.
\end{tcolorbox}
